% Preface
% ~4 pages

\chapter*{서문}
\addcontentsline{toc}{chapter}{서문}
\markboth{서문}{서문}

\section*{동기}

우리는 세상을 물건들의 집합으로 경험한다. 컵, 나무, 얼굴---각각은 독립적이고 경계가 뚜렷한 개체로 인식된다. 그러나 이 경험 이면에는 풀리지 않은 근본적 질문이 있다: 연속적인 감각 흐름으로부터 어떻게 이산적 대상이 형성되는가?

이 책은 이 질문에 대한 하나의 수학적 답을 제시한다. \textbf{Soft Cognitive Cohesion}(SCC) 이론은 물건을 출발점으로 삼지 않는다. 대신 \emph{응집}---관계적 자기-지지의 연속적 장---을 원초적 구조로 삼고, 물건은 이 장이 네 가지 구조적 조건을 충분히 만족할 때 사후적으로 읽어낼 수 있는 안정화된 패턴으로 취급한다. 이 전환은 단순한 일반화가 아니다. 전통적 객체 중심 접근이 ``너무 늦게 시작하는'' 문제---개별화가 이미 완료된 상태를 전제하는---를 정면으로 해결하는 패러다임 전환이다.

SCC는 변분적 이론이다. 응집 장은 네 항의 에너지 범함수를 최소화하며, 각 항은 독립적 구조적 요구---닫힘, 분리, 형태학, 시간적 계승---를 반영한다. 이 에너지 구조 위에서 위상 전이, 기울기 흐름 수렴, $\Gamma$-수렴 등 48개의 정리가 증명되었고, 174개의 테스트를 통과하는 Python 구현이 존재한다.


\section*{이 책의 독자}

이 책은 세 부류의 독자를 염두에 두고 쓰였다. 각 독자에게 가장 효율적인 읽기 경로를 제안한다.

\paragraph{수학자.}
변분법, 위상 장 이론, 그래프 스펙트럼 이론에 관심이 있는 독자라면, Chapter~1에서 이론의 동기를 파악한 후 곧바로 Part~II(Chapter~4--9)의 형식 이론으로 진행할 수 있다. Part~II는 공리 체계, 에너지 범함수의 정확한 정의, 위상 전이 정리(T8), 기울기 흐름 수렴(T14), $\Gamma$-수렴(T11) 등 핵심 정리의 완전한 증명을 포함한다. 부록~A는 필요한 배경(Łojasiewicz 부등식, 스펙트럼 그래프 이론)을 제공한다.

SCC의 수학적 독자성은 에너지의 자기-참조적 구조에 있다: 분리 항 $\mathcal{E}_{\mathrm{sep}}$에서 장이 ``무엇이 구별인가''를 정의하면서 동시에 그 기준에 의해 평가받는다. 이는 표준 Allen-Cahn이나 Ginzburg-Landau 이론에서 나타나지 않는 구조이다.

\paragraph{인지과학자.}
지각적 조직화, 전주의적 처리, 의식의 통합 이론에 관심이 있는 독자라면, Part~I(Chapter~1--3)에서 이론의 철학적 동기와 핵심 개념을 충분히 파악할 수 있다. 특히 Chapter~3은 응집 장의 의미, 전-객체적 인식의 수학적 모델링, 게슈탈트 원리와의 연결을 다룬다. Chapter~7(시간적 지속성)은 발생적 정체성(genidentity) 개념을 형식화하며, Chapter~10--14는 인지과학적 해석과 실험적 예측을 제시한다.

수식에 익숙하지 않은 독자도 Part~I만으로 이론의 핵심 아이디어를 이해할 수 있도록 설계되었다. 각 수학적 개념은 직관적 비유와 구체적 예제로 먼저 소개된다.

\paragraph{ML/AI 연구자.}
표현 학습, 물체 발견(object discovery), 비지도 분할에 관심이 있는 독자라면, Chapter~1에서 동기를 파악한 후 Part~II의 에너지 구조(Chapter~4--5)와 Part~IV의 알고리즘 구현(Chapter~12--14)으로 진행할 수 있다. SCC의 에너지 범함수는 미분 가능하며, 기울기 기반 최적화로 효율적으로 풀 수 있다. Python 패키지 \texttt{scc}는 그래프, 연산자, 에너지, 최적화, 진단을 모듈화하여 제공한다.

% [Figure placeholder: reading paths diagram]
\begin{figure}[t]
\centering
\fbox{\parbox{0.85\textwidth}{\centering\vspace{3.5cm}
\textbf{[그림]} 세 가지 읽기 경로.\\
수학자: Ch1 $\to$ Ch4--9 $\to$ 부록A.\\
인지과학자: Ch1--3 $\to$ Ch7 $\to$ Ch10--14.\\
ML 연구자: Ch1 $\to$ Ch4--5 $\to$ Ch12--14.
\vspace{0.5cm}}}
\caption{청중별 추천 읽기 경로. 실선은 필수 경로, 점선은 선택적 보충을 나타낸다. 모든 경로는 Chapter~1에서 시작한다.}
\label{fig:reading-paths}
\end{figure}


\section*{사전 요구사항}

이 책은 다음의 수학적 배경을 전제한다:
\begin{itemize}
    \item \textbf{선형대수}: 고유값/고유벡터, 양정치 행렬, 스펙트럼 분해. Part~II의 스펙트럼 그래프 이론에 필수적이다.
    \item \textbf{실분석}: 수렴, 콤팩트성, 연속성. 에너지 최소화자의 존재성 증명에 사용된다.
    \item \textbf{그래프 이론}: 인접 행렬, 라플라시안, 연결 성분. SCC의 관계적 지지 공간은 그래프로 모델링된다.
\end{itemize}

PDE 배경(Allen-Cahn 방정식, 기울기 흐름)은 도움이 되지만 필수는 아니다---부록에서 필요한 내용을 제공한다. Python 기초 지식은 Part~IV의 구현 장을 읽는 데 필요하다.


\section*{책의 구조}

이 책은 네 개의 Part로 구성된다.

\textbf{Part~I: 기초} (Chapter~1--3)는 이론의 동기, 핵심 개념, 철학적 기반을 다룬다. 왜 전통적 접근이 한계에 부딪히는지, 응집 장이라는 대안이 무엇을 제공하는지, 그리고 proto-cohesion 진단 벡터가 어떻게 형성의 질을 평가하는지를 소개한다.

\textbf{Part~II: 형식 이론} (Chapter~4--9)은 SCC의 수학적 핵심이다. 형식 우주와 공리 체계(Chapter~4), 에너지 범함수와 변분 구조(Chapter~5), 핵심 정리들---존재성(T1), 위상 전이(T8), 수렴(T14), $\Gamma$-수렴(T11)---의 증명(Chapter~6--9)을 포함한다.

\textbf{Part~III: 시간적 구조와 다중 형성} (Chapter~10--11)은 시간적 전달과 지속성의 이론(Chapter~10), 그리고 여러 형성이 공존하는 다중 형성 체계---K-필드 아키텍처, 운동학적 장벽, 결합 매개변수 $\Lambda_{\mathrm{coupling}}$---(Chapter~11)을 다룬다.

\textbf{Part~IV: 구현과 응용} (Chapter~12--15)은 Python 패키지 \texttt{scc}의 구조(Chapter~12), 수치 실험과 검증(Chapter~13), 인지과학적 연결과 응용(Chapter~14), 완전한 정리 레지스트리(Chapter~15)를 제시한다.


\section*{이 책을 읽는 방법}

\paragraph{순차적 경로.}
처음부터 끝까지 순서대로 읽는 것이 가장 완전한 경험을 제공한다. Part~I에서 직관을 쌓고, Part~II에서 수학을 전개하며, Part~III에서 시간과 다중 형성으로 확장하고, Part~IV에서 구현과 응용으로 마무리한다.

\paragraph{전문화된 경로.}
시간이 제한되거나 특정 관심사가 있는 독자는 위의 ``이 책의 독자'' 섹션에서 제안한 경로를 따를 수 있다.

\paragraph{Part 간 의존성.}
Part~I은 독립적으로 읽을 수 있다. Part~II는 Part~I의 개념적 이해를 전제한다. Part~III는 Part~II의 에너지 구조를 전제한다. Part~IV는 Part~II의 주요 정의를 참조하지만, 증명의 세부사항 없이도 읽을 수 있다.

\paragraph{표기법.}
이 책은 Canonical Specification v2.1의 표기 규약을 따른다. 다음 ``표기법'' 장에서 주요 기호와 규약을 정리한다. 순서쌍 합(ordered-pair summation) 규약---모든 $\sum_{x,y \in X_t}$는 순서쌍을 범위로 함---은 위상 전이 정리(T8)와 Hessian 분석에서 핵심적이므로 반드시 숙지해야 한다.

\paragraph{코드 예제.}
본문에 포함된 Python 코드는 개념을 구체화하기 위한 것이다. 완전한 구현은 \texttt{scc/} 패키지에서 찾을 수 있다. 코드를 실행하려면 Python~3.9 이상, NumPy, SciPy가 필요하다.

\bigskip
\noindent
이 책이 연속적 인식과 이산적 대상 사이의 간극을 탐구하는 데 유용한 안내서가 되기를 바란다.

\vspace{1cm}
\begin{flushright}
오재홍\\
2026년 4월
\end{flushright}
